% Venue-neutral body. Included by arxiv.tex / workshop.tex.
% Do not put \documentclass, packages, title or bibliography here.
% The abstract and keywords live in the wrapper (arxiv.tex / workshop.tex),
% not here: every venue formats them differently and some classes require
% them before \maketitle. They were duplicated in both places until 2026-08-04,
% which compiled without error and printed the abstract twice.

\section{Introduction}

When a disaster is large enough, people go toward it. This is among the
oldest findings in disaster sociology. \citet{fritz1957} named it --- \emph{``the informal, spontaneous movement of people, messages,
and supplies toward the disaster area''} --- and catalogued five kinds
of person who arrive: the returnees, the anxious, the helpers, the
curious, and the exploiters. Their report is subtitled \emph{A Problem
in Social Control}, and the observation it opens with has not dated:
convergence \emph{``brings needed aid to many victims, but at the same
time the resultant congestion makes organization and control of the
rescue and relief efforts more difficult.''} \citet{kendra2002}, studying the response to the World Trade Center attack, added a
sixth and reframed the central question. Access to what they call the
\emph{response milieu} is not granted; it is negotiated. A volunteer's
admission depends on whether the people already inside can afford them.

The mechanism the field has developed for that negotiation is
\textbf{credentialing}. Volunteer reception centres, badge systems,
affiliation with a recognised organisation: these convert an unknown
arrival into an accountable one. They work. At the World Trade Center,
identification requirements \emph{``evolved and intensified on almost a
twice-daily basis''} as the response matured.

Practitioner doctrine is more explicit still, and defines our population
for us. \citeauthor{fema_msv}'s guidance on managing spontaneous volunteers \citeyearpar{fema_msv} describes
unaffiliated volunteers as those who are \emph{``no part of a recognized
voluntary agency''}, who \emph{``often have no formal training in
emergency response''}, and who --- the phrase is theirs --- are
\emph{``not officially invited to become involved.''} It adopts the
academic taxonomy directly, noting that \emph{``researchers have
identified six different groups of people that tend to converge''}, and
identifies the operational task: \emph{``the helpers must be identified
from among the larger population of convergent individuals.''}

The doctrine offers two mechanisms for doing that identification, and
both presuppose an institution. The first is the Volunteer Reception
Centre, which processes arrivals --- and must therefore have been
established. The second, and the one the guidance treats as primary, is
prevention:

\begin{quote}
``Turn spontaneous unaffiliated volunteers into affiliated ones
\textbf{before a disaster occurs.} People who make a pre-disaster
decision to become disaster volunteers and take training to prepare
themselves will NOT become spontaneous, unaffiliated volunteers after a
disaster.''
\end{quote}

The same materials pose, as a training exercise, the question of how a
community can \emph{``keep your community members from
self-deploying''}, and describe organised national service volunteers
approvingly as members who \emph{``never self-deploy, but wait to be
called.''}

We take this seriously rather than dismissively: pre-affiliation plainly
works, and a trained volunteer who arrives when called is better for
everyone than one who does not. But both mechanisms require an
institution --- one that reached the volunteer months beforehand, or one
that has arrived and opened a reception centre. Neither covers the
person who decides at two in the morning, on the night, that they have a
boat. Doctrine's answer to that person is that they should not exist.
They do exist, in thousands, and the events that produce them are the
events where the institution is least able to arrive.

\citet{kendra2002} state the resulting tension directly, summarising
Weick and Perrow on decentralised coordination: effective
decentralisation presupposes prior socialisation into an organisation's
norms, \emph{``yet the volunteers who appear to assist in the emergency
response are, virtually by definition, strangers to the response
milieu.''} The literature names the problem and, so far as we have
found, proposes no mechanism for the interval before command exists.

Hurricane Harvey makes the interval concrete. In August 2017 flooding in
the Greater Houston area affected an estimated 30,000--40,000 homes, and
city officials and FEMA publicly asked citizens with boats to help reach
people trapped inside them \citep{smith2018}. Thousands came. They
organised through Zello, a push-to-talk application, along with Facebook
groups, NextDoor, ad hoc Google spreadsheets, and applications written
during the event itself. Interviewing twenty of these volunteers, Smith
et al.~found at least twenty distinct groups with, between them,
\emph{``no way to seamlessly share information and coordinate
activity.''} Affiliation was nominal: membership of the loose ``Cajun
Navy'' was \emph{``fleeting''}, leadership \emph{``fluid and
dispersed''}, boundaries unclear.

What went wrong is documented in the volunteers' own words. One
dispatcher described the coordination problem exactly:

\begin{quote}
``it became overwhelming just trying to keep track of who was going out
who was coming back who was out on a boat, \textbf{can we account for
everyone}, and who's in their house and it was a mess.'' --- Gary, in
\citet{smith2018}
\end{quote}

The gap this paper addresses is not one we inferred. It was stated by a
person who was in it.

\citeauthor{smith2018} name three coordination failures --- incomplete feedback
loops, unclear prioritisation, and communication overload --- and close
by calling for \emph{``the design of intuitive systems that can quickly
be mastered by the novice social media user.''} This paper is an attempt
at one narrow part of that.

Two details from that study bear directly on the design that follows.
First, the accountability function was being performed --- badly --- by
\emph{families}: rescuers reported phones ``constantly blowing up'' with
relatives asking \emph{``are you okay, are you okay, are you okay.''}
Somebody was always going to ask the question; the question simply had
nowhere to be answered. Second, volunteers' phones were destroyed by
rain and floodwater in numbers large enough that many replaced them
afterwards --- which is direct evidence that silence and danger are
genuinely different things, and that any system escalating on silence
will produce false alarms.

This paper reports the design of \textbf{DiresQ}, a system that attempts
accountability for such volunteers without a coordinator. Its central
move is to escalate on \textbf{silence} rather than on supervision: a
responder states an expected arrival time and checks in periodically,
and if they stop, the system files a report about them automatically.
Nobody has to notice.

We are explicit about what this paper is not. DiresQ has never been used
in a real disaster. It has no users, no deployment beyond a public
demonstration instance, and no evaluation. We report a design, the
reasoning behind five specific trade-offs, and the failure modes we
encountered building it --- several of which we consider more
transferable than the system itself. Readers looking for evidence that
this approach works will not find it here, and we would regard any such
reading as a misuse of the paper.

\section{Background and related work}

\subsection{Convergence}

People moving toward a disaster is one of the field's oldest documented
behaviours. \citet{fritz1957} named it \emph{convergence
behavior} --- \emph{``the informal, spontaneous movement of people,
messages, and supplies toward the disaster area''} --- and observed that
it \emph{``brings needed aid to many victims, but at the same time the
resultant congestion makes organization and control of the rescue and
relief efforts more difficult.''} They separate movement \emph{``toward
the struck area from the outside --- external convergence''} from
\emph{``movement toward specific points within a given disaster-related
area or zone --- internal convergence''}, distinguish three forms
(personal, informational, materiel), and catalogue five types of
personal converger: the returnees, the anxious, the helpers, the
curious, and the exploiters.

It is worth noting that the founding treatment frames convergence as a
control problem --- its subtitle is \emph{A Problem in Social Control}
--- and that this framing has been remarkably durable. The paper's own
position is stated in §5.

\citet{kendra2002}, studying the response to the World Trade
Center attack across more than 750 collective hours of field
observation, added a sixth --- supporters or fans --- and, more
importantly for us, reframed the question. Access to what they term the
\emph{response milieu} is not a status a volunteer holds but one they
negotiate.

The taxonomy has passed into practice. \citeauthor{fema_msv}'s guidance \citeyearpar{fema_msv} on managing
spontaneous volunteers reproduces it directly, noting that
\emph{``researchers have identified six different groups of people that
tend to converge''}, and derives an operational task from it:
\emph{``the helpers must be identified from among the larger population
of convergent individuals.''} Theory and doctrine here are a single
lineage rather than two literatures, which is why we cite them together.

DiresQ's users are helpers, converging externally and then internally.

\subsection{Legitimacy, credentialing, and what both presuppose}

\citeauthor{kendra2002}'s \citeyearpar{kendra2002} central finding concerns who gets in. (We cite
their Disaster Research Center preliminary paper throughout, because it
is the text we read; the 2003 book chapter develops the same fieldwork
and is listed in the references for completeness \nocite{kendra2003}.) At the World Trade
Center, identification requirements \emph{``evolved and intensified on
almost a twice-daily basis''}, and volunteers became \emph{``another
group that needed to be accounted for, and therefore potentially a
distraction that outweighed their utility.''} The volunteers who
succeeded in gaining access were those who \emph{``were able to work
with minimal supervision by official emergency workers''} --- whose
incorporation \emph{``required little or no effort on the part of
emergency managers.''}

We take the scarce resource in that account to be emergency-manager
attention, and it is the frame we design against: a volunteer who
accounts for themselves spends none of it.

Doctrine's mechanisms for conferring legitimacy are the Volunteer
Reception Centre and, prior to any event, pre-affiliation ---
\emph{``turn spontaneous unaffiliated volunteers into affiliated ones
before a disaster occurs.''} Both presuppose an institution: one that
has arrived and opened a centre, or one that reached the volunteer
months earlier. We develop the consequence in §1 and do not repeat it
here.

\subsection{Digital volunteers, and why they need no accounting for}

\citet{starbird2011} provide the foundational study of volunteers
self-organising through technology after a disaster, examining 292,928
tweets and interviewing nineteen of the ``voluntweeters'' who emerged
after the 2010 Haiti earthquake. They read the phenomenon through Kreps
and Bosworth's structural theory, finding a resource → activity → task →
domain progression in which the \emph{resource} --- Twitter, and the
individual capacities it made usable --- is what lets a stranger begin.

The distinction that matters for this paper is one their scope makes for
us: their volunteers are remote by definition. The population is people
helping from other continents. No participant is in the hazard, so no
participant needs accounting for, and no mechanism for doing so appears
in the paper. This is not a gap in their work. It is the boundary of
ours.

Our question is what changes when that same self-organising sequence
runs among people who are physically inside the hazard. Accountability
stops being optional and there is still nobody to provide it.

\subsection{Citizen-led response in practice}

\citet{smith2018} interviewed twenty participants in the citizen-led
rescues during Hurricane Harvey and found three roles --- rescuer,
dispatcher, information compiler --- distributed across at least twenty
distinct groups with \emph{``no way to seamlessly share information and
coordinate activity.''} They report three coordination failures:
incomplete feedback loops, unclear prioritisation, and communication
overload.

Incomplete feedback loops is our problem under another name. Their
account of it is mechanical rather than motivational: rescue platforms
depended on someone manually marking a case closed, and in practice the
case number was forgotten, the phone was destroyed, or a different boat
reached the address first. The loop stayed open not because anyone
neglected it but because closing it required an action at exactly the
moment nobody had attention to spare.

We note two of their findings as constraints on our design rather than
support for it. First, their participants wanted \emph{more}
coordination: one dispatcher argued that \emph{``people that are trained
to do that need to be in charge of prioritizing calls and assigning.''}
Second, over-convergence was itself a failure mode --- boats were turned
away from neighbourhoods that already had too many. We return to both in
§5.

\subsection{Systems for crisis crowdsourcing}

\citet{liu2014} offers the field's most developed design framework,
organising a crowdsourcing system around why / who / what / when / where
/ how, and around the social, technological, organisational and policy
interfaces that manage the articulation work of coordinating across
them. We note that two of those four interfaces presuppose an
organisation conducting the effort, and that the trajectory the
framework describes runs toward formalisation and integration
\emph{``into official products and services.''} Our position is upstream
of that trajectory rather than opposed to it.

\citet{auferbauer2017} describe \emph{crowdtasking}, presented
explicitly as \emph{``a centralized form of crowdsourcing for crisis and
disaster management''}, with a prototype and a first field trial. It is
the clearest statement of the opposite design position to ours, and we
cite it as such: where crowdtasking assumes an organisation able to task
a registered crowd, we assume neither the organisation nor the
registration.

\citet{kankanamge2019} provide a systematic review of volunteer
crowdsourcing in disaster risk reduction. We have not obtained the full
text and therefore make no claim about its findings; it is listed here
because a reader working in this area will expect to see it and should
know we are aware of it.

Technical work on the classification side is considerably more advanced
than ours. \citet{zhou2022} fine-tune BERT variants over a
hand-labelled corpus of Hurricane Harvey tweets to identify rescue
requests, and substantially outperform the earlier baselines they
compare against. Our own priority suggestion (§3.3) is a naive Bayes
classifier over a small hand-labelled corpus, and is not competitive
with this work, nor intended to be: it exists to run in a browser with
no network, which is a different constraint rather than a weaker attempt
at the same one.

\subsection{Where this leaves the gap}

Convergence is thoroughly theorised. Legitimacy is theorised and
operationalised. Digital volunteering is well studied among remote
participants. Crowdsourcing systems have a design framework and at least
one field-trialled centralised implementation. Classification of rescue
requests is a mature technical problem.

What we have not found, in the literature or the doctrine, is a
mechanism for accounting for physically converging helpers during the
interval before an institution exists to do it. Doctrine's answer is
that those people should have been affiliated beforehand. The academic
literature describes their arrival as a management problem. Neither
addresses the person already in the water.

That interval is this paper's subject, and the design in §3 is one
attempt at it.

\section{Design}

DiresQ is a server-rendered web application: Flask, SQLite, no client
framework, no account approval step, and no dispatcher role.
Thirty-three routes over six tables. The stack is deliberately
unremarkable, and we describe it only to establish that nothing in what
follows depends on infrastructure a volunteer organisation would not
already have. The contribution is not the stack.

Anyone can sign up. There is no verification of identity, no vetting,
and no administrator who admits people. This is not an oversight we
intend to fix; it is the condition the system is designed for. A tool
that requires somebody to approve you has reintroduced the coordinator
whose absence is the entire problem.

\subsection{The mechanism}

A report is a place that needs help. A responder who decides to go there
joins the report and, in doing so, states an expected time of arrival.
The ETA is bounded between five and 240 minutes and warns above 120 ---
not because longer journeys are invalid, but because an unbounded ETA
makes the subsequent arithmetic meaningless. From that point the
responder is expected to check in; the interval defaults to thirty
minutes.

If a responder stops checking in, they are marked overdue. Fifteen
minutes past that threshold, the system files a report \emph{about them}
--- their last known location becomes an incident in the same feed
everyone else is reading. No human decides this. No human is asked to
notice.

This is the whole mechanism, and its only novel property is what
triggers it. Existing accountability systems escalate when a coordinator
observes that somebody has not reported in. We escalate on the
observation itself being absent. Silence is available in circumstances
where a supervisor is not.

\subsection{Five decisions, and what each one cost}

The design is small enough that the interesting content is not the
architecture but the trade-offs. We describe five, each of which we got
wrong first.

\textbf{Silence as the trigger.} The alternative designs all require
somebody to be watching: a dispatcher view, an alert queue, a supervisor
role. Each of them works, and each of them assumes the thing we cannot
assume. Escalating on absence needs nobody. The cost is that absence is
ambiguous --- a responder who has stopped answering may be in trouble,
or may have a dead phone, or may have gone home without saying so. The
system cannot tell these apart and does not claim to. It reports
\emph{contact lost}, never \emph{safe} or \emph{unsafe}, and the
interface language was rewritten twice to stop implying otherwise.

\textbf{No in-process timer.} Our first design ran the silence check on
a background thread inside the application. We removed it. A thread that
dies takes the alarm with it and leaves the interface showing green, and
a green screen is read as a positive result rather than an absent one.
The check now runs on read: whenever anyone loads the accountability
board, the sweep runs first, rate-limited to once every thirty seconds.
An \emph{external} scheduler is supported and optional
(\texttt{flask\ -\/-app\ app\ sweep} from cron or Task Scheduler) ---
deliberately external, because an external scheduler that fails is at
least visible to the machine running it, which an in-process timer is
not.

This trades one dependency for another. A timer that might die becomes a
check that depends on being watched. We consider that the better failure
--- an unwatched board is a situation where nobody is relying on the
result --- but the substantive move is that we made the dependency
\textbf{visible} rather than documenting it. The board displays how long
since the last sweep and turns amber past five minutes. On a board
somebody is actually watching, the number always reads a few seconds,
because the watching is what runs it. A limitation you can see on screen
is a different claim from one the user has to be told about.

\textbf{Status and contact are different facts.} What a responder last
told us is one thing; whether they are still answering is another. We
conflated them, and the result was a report page displaying ``on scene''
for somebody the accountability board had forty-five minutes overdue.
Anyone opening that report to decide whether the address needed more
help was counting a person who had gone silent as help present. Both
facts are now shown, and neither is inferred from the other. Somebody
who has explicitly cleared is not chased, because going home is not
going quiet.

The general form of this: in a system whose purpose is knowing where
people are, \emph{self-reported state} and \emph{liveness} must never be
collapsed into one field, however tempting the simplification is at the
schema level.

\textbf{Refusing to cache claims about other people.} The offline layer
keeps the things you committed to --- your own assignments, a queued
report you filed without signal --- and refuses to keep the report feed
or the accountability board. A cached feed is a list of who needed help
twenty minutes ago. Acting on it sends somebody to an address that has
been cleared, and the person reading it has no way to know the
difference between stale and current. A test fails if either page is
added to the service worker's shell list.

The generalisable claim, and the one we would defend:
\textbf{correct-when-written and correct-when-read are different
guarantees}, and disaster data has an unusually short half-life between
them. Offline-first design literature tends to treat availability as an
unalloyed good. For a class of data it is not, and the distinction is
not about staleness tolerance but about whether a stale answer is
\emph{actionable} --- whether a user can be sent somewhere by it.

\textbf{Not storing triage answers.} The system includes a START triage
helper that orders which reports get attention first. The answers are
health observations about a person who did not consent and is likely in
no position to. They are used to compute an ordering and then discarded;
nothing about a casualty's breathing or perfusion is written to the
database. The cost is that the ordering cannot be audited after the
fact, which is a real loss. We accepted it.

\subsection{Priority suggestion}

A report can be given a suggested priority by a multinomial naive Bayes
classifier over a severity lexicon, trained on a small hand-labelled
corpus and validated by leave-one-out cross-validation with a floor
asserted in continuous integration.

We describe it in one paragraph deliberately. It is the least
interesting component and the one most likely to be mistaken for the
contribution. It suggests an ordering for human attention; it does not
predict outcomes, assess medical severity, or make any claim that
survives the person reading the report disagreeing with it. A model
small enough to ship as a frozen table of word counts is also small
enough to run in the browser, which is why priority suggestion works
with no signal while the report feed deliberately does not.

\subsection{What the design does not attempt}

It does not summon help. No dispatcher reads it, it does not call
emergency services, and the disclaimer on every page is load-bearing
rather than legal decoration. It does not verify that anybody is who
they say they are. It does not know where anyone actually is ---
location is self-reported throughout. It makes no claim that a responder
is safe; only that we have or have not heard from them.

These are stated here rather than only in §5 because a reader who
reaches the limitations section still holding the wrong model of the
system has already misread the design.

\section{What building it taught us}

We report these as observations from construction, not as validated
findings. Each began as a specific bug or reversal and generalised
afterwards.

\subsection{Correct-when-written and correct-when-read are different guarantees}

The offline layer was originally designed to cache what a user would
want if their signal dropped, which is nearly everything. We now cache
almost nothing.

The rule we arrived at: \textbf{do not store anything that stops being
true when it is written.} Your own commitments --- the assignment you
accepted, a report you filed with no signal --- remain true offline,
because they are statements about you. The report feed and the
accountability board are statements about other people, and both decay.
A cached feed is a list of who needed help twenty minutes ago; a
responder acting on it is dispatched to an address that has since been
cleared, with no way to distinguish that from a current one.

Offline-first design generally treats availability as an unalloyed good,
with staleness managed by tolerance windows and revalidation. We think
the useful distinction is not how stale the data is but whether a stale
answer is \emph{actionable} --- whether a user can be sent somewhere by
it. Data that can move a body should not be served from cache at all. A
test fails if either page is added to the service worker's shell list,
because the rule is easy to forget and the failure is silent.

\subsection{An alarm that dies quietly is worse than no alarm}

The silence check was first designed as a background thread inside the
application. We removed it. A timer that dies takes the alarm with it
and leaves the interface showing green, and a green screen is read as a
positive result rather than an absent one.

The check now runs on read --- whenever anyone loads the accountability
board, rate-limited to once every thirty seconds. An external scheduler
is supported (\texttt{flask\ -\/-app\ app\ sweep}) and deliberately
optional: an external scheduler that fails is at least visible to the
machine running it, which an in-process timer is not.

This does not remove the dependency, it relocates it. The system now
works while somebody is looking. We regard that as the better failure
mode, because a board nobody is watching is a situation where nobody is
relying on the answer --- but the claim needs stating rather than
assuming.

\subsection{A limitation on screen is a different claim from a limitation in a file}

The consequence of §4.2 is that the guarantee is conditional, and our
first response was to write that down in a limitations document. That is
the standard move and we now think it is insufficient.

The board displays how long since the last sweep --- \emph{checked 2s
ago} beside the live indicator, amber past five minutes, which is ten
minutes before the fifteen-minute escalation it drives. On a board
somebody is watching it always reads a few seconds, because the watching
is what runs it. If it stops moving, that is visible too.

The general form: \textbf{emergency software should surface its own
liveness.} A user deciding whether to trust a screen needs to know
whether the thing behind it ran, and that is a different question from
whether the data looks reasonable. We would extend this beyond our own
system: any interface that reports on a periodic check should display
when the check last completed, not merely its result.

\subsection{The pages built to be watched were the ones that broke under watching}

The board, the map and the feed all run the silence check before
answering, and the check writes. SQLite's default journal mode lets one
writer lock out every reader. With several people watching a board that
refreshes every three seconds, a write could wait past its timeout and
fail the entire response.

The pages whose purpose is to be watched continuously during an
emergency were precisely the pages that failed when watched
continuously. The fix was unremarkable --- write-ahead logging,
rate-limiting the sweep, catching the failure rather than letting it
become an error page. The lesson is not about SQLite. It is that
\textbf{making a read-only page write is a change of kind, not of
degree}, and in this system it converted the most-loaded pages into the
most fragile ones.

One deliberate detail: a failed sweep does not record itself as having
run. The timestamp goes stale and the board reports that in amber, which
is true, because a check we could not record is not a check we can
claim.

\subsection{Prose about a system lags the system}

Our test suite verifies every count in our own documentation --- routes,
tables, tests, lines --- and fails when the prose and the repository
disagree. It has caught wrong numbers repeatedly, including several we
wrote ourselves.

It has never once caught the more common error, which is a
\emph{sentence} that describes a previous version of a behaviour. Four
times a change landed and the paragraph explaining it stayed a revision
behind. Numbers are checkable and were checked; claims are not, and were
not.

We do not have a solution and offer this as an observation: automated
documentation checks create a real assurance about the class of
statement they can verify and no assurance whatever about the class they
cannot, and the second class is where the misleading statements live.

\section{Limitations}

The system has eighteen documented limitations, maintained as a
first-class document rather than an appendix. We summarise the ones that
most constrain what this paper claims.

\textbf{It has never been used in a real disaster.} Everything here is
reasoned from published accounts of Harvey, Kathmandu and Mexico City,
and from published triage protocol. None of it has been tested by
somebody standing in water at two in the morning. We believe the
reasoning is sound; that is not the same as knowing it works, and no
part of this paper should be read as evidence that it does.

\textbf{The switch depends on being run.} Described in §4.2. Neither the
read-trigger nor an external scheduler is guaranteed. A deployment
nobody looks at and nobody configures performs no checks.

\textbf{Overdue measures contact, not safety.} The system reports that
it has not heard from somebody. It cannot distinguish danger from a dead
battery, and it does not try. The interface language was revised twice
to stop implying otherwise, and this is the limitation most likely to be
misread by a user in a hurry.

\textbf{No identity verification.} Anyone can register. This is a design
condition rather than a defect --- a system requiring approval has
reintroduced the coordinator whose absence is the problem --- but it
means the system cannot distinguish a responder from someone claiming to
be one.

\textbf{Location is self-reported} throughout. The map shows where
people said they were.

\textbf{One cautious responder can hold a report open}, which biases the
system toward over-reporting need. We prefer that direction and note
that we chose it.

\textbf{Lockouts live in memory} and do not survive a restart.

\textbf{The volunteers we claim to serve asked for the opposite.} Smith
et al.~(2018) interviewed twenty Harvey rescuers and dispatchers, and
the structural complaint that came back was a wish for \emph{more}
coordination, not less. One dispatcher, on being unable to prioritise
calls: \emph{``people that are trained to do that need to be in charge
of prioritizing calls and assigning.''} Another described sending boats
away because too many had converged on one neighbourhood.

This is the strongest available argument against our design, and it
comes from the population the design is for. We do not think it defeats
the argument --- the wish for a trained coordinator does not summon one,
and the interval before one exists is exactly our subject --- but a
reader should weigh it. A system that makes self-deployment easier and
better-accounted-for may also make it more attractive, and \citeauthor{smith2018} document a response that suffered from too many boats as well as too
little information.

\subsection{A position rather than a finding}

The literature we build on treats converging volunteers as a population
to be managed: \citet{kendra2002} document access as negotiated
legitimacy, with credentialing, liability and security as the operative
concerns, and volunteers as a potential distraction that must outweigh
its own cost. Our design treats the volunteer as somebody owed an
accounting.

These are different moral starting points, and the literature does not
share ours. We think ours is defensible --- a person who walks into a
hazard to help is owed something regardless of whether an institution
has authorised them --- but it is a position, not a result, and a reader
is entitled to reject it and evaluate the design on the field's own
terms instead.

\section{Future work}

The obvious next step is the one we have not taken: contact with people
who do this work. Five structured interviews with volunteer-response
coordinators --- people who have run a reception centre, or turned
volunteers away, or gone in unaffiliated themselves --- would test the
assumption the whole design rests on, which is that the interval before
credentialing exists is a real operational gap rather than an artefact
of how we read the literature.

Beyond that, a simulated-event walkthrough is the tractable route to
evaluation without waiting for a disaster; the crisis informatics
community has precedent for this. It would let us observe whether the
silence mechanism produces escalations at useful times, or whether it
mostly produces false alarms from people whose phones died --- a
distinction we currently cannot make.

Two smaller questions we would want answered. Whether the fifteen-minute
escalation delay is anywhere near right; we chose it by reasoning, not
by measurement, and it is the kind of parameter that ought to come from
data. And whether making system liveness visible on screen --- the sweep
timestamp described in §4.3 --- actually changes how much a user trusts
what they are reading, or whether it is a designer's satisfaction that
no user notices.

We are aware that all of this describes work requiring access to a
professional community that two secondary-school students do not have.
We would welcome collaboration, and we mention it here rather than in an
acknowledgement because it is a limitation on the research, not a
courtesy.

% ---------------------------------------------------------------------------
% STATEMENT ON THE USE OF GENERATIVE AI
%
% Two versions below. Use ONE. Delete the other.
%
% Version A is accurate for the text as it currently stands.
% Version B becomes accurate only once the authors have rewritten the prose in
% their own words. Do not use B before that is true.
%
% Note: SocArXiv's AI Policy (8 March 2026) lists "generating text which is used
% verbatim (including whole paragraphs and sections)" as an unacceptable use.
% Version A therefore describes a use that SocArXiv does not accept. It is here
% so the statement is honest, not so that it passes.
% ---------------------------------------------------------------------------

% ---- VERSION A: accurate for the current text ----
\section*{Statement on the use of generative AI}

The text of this paper was drafted by a large language model (Claude,
Anthropic) from sources and an argument structure specified by the authors.
The authors set the research question, chose the sources, made every design
decision reported in Section~3, and reviewed and revised the result. Sections
of the drafted text appear substantially as generated. Three cited works
(Liu 2014; Auferbauer and Tellio\u{g}lu 2017; Kankanamge et al. 2019) are
referenced only for positions stated in their abstracts, which is noted where
they appear.

% ---- VERSION B: use only after the authors rewrite the prose ----
% \section*{Statement on the use of generative AI}
%
% A large language model (Claude, Anthropic) was used for literature search,
% source retrieval, and organisation of material during the preparation of this
% paper. All prose is the authors' own. The authors verified every source cited,
% including confirming that quoted passages appear in the works attributed to
% them, and made every design decision reported in Section~3.
